\begin{figure*}[t]
\centering
\begin{tcolorbox}[
  colback=gray!5,colframe=gray!50,boxrule=0.4pt,
  fontupper=\small\ttfamily,left=6pt,right=6pt
]
You are reading a long document one page at a time to answer a question.
Pages arrive in retrieval-ranked order, not document order. You have no
memory of earlier turns beyond NOTES. A later turn will not see the current
page, so record any relevant evidence it contains.\\[0.5em]

Think step by step. Then output these markers on separate lines:\\
REASONING: <your assessment of the page and notes>\\
RECORD: <evidence from this page, or exactly ``none''>\\
ANSWER PAGE: <true or false>\\
DECISION: <CONTINUE or STOP>\\[0.5em]

Record only what the CURRENT PAGE adds. Preserve names, numbers, units,
table headers, and the details needed to interpret visual evidence. For a
count or list, identify matching items on this page and where each appears;
end the record with COUNT ON THIS PAGE: <number> (<what was counted>).
Do not copy evidence already in NOTES.\\[0.5em]

Mark ANSWER PAGE: true whenever the answer may depend on this page,
including when it supplies only part of a multi-page answer or you are
uncertain of its relevance. Flagged pages will be shown to the synthesizer
in full.\\[0.5em]

Before the applicable reading floor, write DECISION: CONTINUE. After that
floor, write DECISION: STOP only when the accumulated evidence appears
sufficient to answer the question completely. A later ranked page may
still contribute to a count or list.\\[0.5em]

QUESTION: \{question\}\\
NOTES: \{reader records from earlier pages\}\\
CURRENT PAGE: \{document page number, parser text, page image\}
\end{tcolorbox}
\caption{Reader prompt structure for iteration 7. The stopping instruction uses the initial floor $k$ or the resumed floor $m$, as appropriate.}
\label{fig:app_reader_prompt}
\end{figure*}